% s-376-p376-m001; original PDF 376
\PoemStart{s-376-p376-m001}{Der blaue See von Lucel.}
\begin{BookVerse}
\VerseLine{0}{Du bist so rein, und deine zarte Bläue}
\VerseLine{1}{Zeigt spiegelnd des Gebirges feinste Züge.}
\VerseLine{1}{In deinem Blick birgt sich kein Fältchen Lüge,}
\VerseLine{0}{Durchsichtig bis zum Grund ist deine Treue.}
\Stanza
\VerseLine{0}{Dein Wesen ist so schamhaft, dass ich scheue}
\VerseLine{1}{Aus dir zu trinken; fürcht’ ich doch, ich füge}
\VerseLine{0}{Dir schon ein Leid zu, wenn ich zur Genüge}
\VerseLine{0}{Mein Aug’ an deinem keuschen Glanz erfreue. —}
\Stanza
\VerseLine{0}{Weh! ärger ward noch nie ein Bild entweiht!}
\VerseLine{0}{Was muss ich sehn! Alt-Englands dürrste Maid —}
\VerseLine{1}{Oh, mehr bedurft’ es nicht, mich zu vertreiben!}
\Stanza
\VerseLine{0}{Im Nymphenweiher solch ein grät’ger Fisch!}
\VerseLine{0}{Oh Geist des Sees! mit einem Binsenwisch}
\VerseLine{1}{Lass mich den frechen Vorwitz trocken reiben!}
\end{BookVerse}
\PoemEnd
